% page 431 du fac-simile (image p-430 du scan)
% bloc 165.1 x 242.0 mm ; marge gauche 36.9 mm
\pgc{15.240mm}{0.000mm}{
\l{\cel{50}423}
\l{}
\l{\sou{parietario.}\cel{2}(\sou{bot.)}\cel{3}Planto ek la familio \textquotedbl{}urticei\textquotedbl{}, qua kres-}
\l{\cel{3}kas sur la muregi. L.\cel{1}\sou{parietaria}.\cel{2}DEFIS.}
\l{}
\l{\sou{parko}.\cel{2}Extensajo ek sulo, cirkondata da muro, greto, e c., qua}
\l{\cel{3}dependas generale de kastelo, de domego, ed uzesas sive por}
\l{\cel{3}konservar vildo, sive por plubeligar rezideyo, sive por esar}
\l{\cel{3}promeneyo publika. - DEFIRS.}
\l{}
\l{\sou{parlamento}.\cel{2}I. Asemblitaro ek personi qui deliberas. - II.}
\l{\cel{3}En Britania : Asemblitaro ek la baroni, prelati e deputiti di}
\l{\cel{3}la urbi,qua helpas la rejo en lua praktiko di la povo suprega.}
\l{\cel{3}- III. Nomo di la asemblitaro, di la \textquotedbl{}chambri\textquotedbl{}, qui praktikas}
\l{\cel{3}la povo suprega en la landi konstitucala. - DEFIRS.}
\l{}
\l{\sou{parlementar}.\cel{2}(\sou{netrans.}) Konferar, negociar, espere di trans-}
\l{\cel{3}akto, koncilio. - DEFIRS.}
\l{}
\l{\sou{parodiar}. (\sou{trans.})\cel{2}Imitar, deformante (lu) burleske, verko}
\l{\cel{3}grava. - DEFIRS.}
\l{}
\l{\sou{paroko}. (\sou{religio).}\cel{2}Sacerdoto qua direktas distrikto ekleziala,}
\l{\cel{3}che la katoliki. - DEFIS.}
\l{}
\l{\sou{parolar}.\cel{2}I. (\sou{netrans.})\cel{2}Uzar la artikulo-linguo. - II. Ex-}
\l{\cel{3}presar sua penso per la artikulo-linguo. - III. (\sou{trans}.)}
\l{\cel{3}Uzar, kom expresilo (ta o ca linguo). - FIS.}
\l{}
\l{\sou{paromo}. Batelo plata, uzata por transportar, sur rivero o lago}
\l{\cel{3}o fluvio, de ca a ta rivo,-homi, animali, veturi, e quan onu}
\l{\cel{3}glitigas alonge kablo tensita inter la du rivi, oblique ye}
\l{\cel{3}la fluo. - F.}
\l{}
\l{\sou{paronimo}.\cel{2}(\sou{gram.})\cel{2}Preske homonimo. DEFIRS.}
\l{}
\l{\sou{paronomiazo}.\cel{2}(\sou{retori}ko).\cel{2}Interproximigo, en frazo, di vorti}
\l{\cel{3}qui pronuncesas inter-same, ma di qui la senci interdiferas.}
\l{}
\l{\sou{parotido}. (\sou{anat.)}\cel{2}Glando salivifera, situita apud la orelo. EFIS}
\l{}
\l{\sou{paroxismo}. (\sou{patol.})\cel{2}Intenseso maxima di aceso. DEFIRS.}
\l{}
\l{\sou{parqueto}.\cel{2}Asemblajo ek lami de ligno harda, polisita, dispozita}
\l{\cel{3}diverse, de qui onu facas la planko-suli eleganta. DEFR.}
\l{}
\l{\sou{parto}. I. La porciono di kozo qua apartenas a singla persono}
\l{\cel{3}kande onu dividas lu inter plura personi, - II. To quo (bona}
\l{\cel{3}o mala) apartenas ad ulu, en kozo en qua li havas interesto.}
\l{\cel{3}- III. Un de la elementi ye qui konsistas loko. - IV. Ele-}
\l{\cel{3}mento ek toto. - V. Elemento qua kontributas ad ensemblo.}
\l{\cel{3}- VI.\cel{1}\sou{(Muziko)}. To quon singla voco havas kantenda - singla}
\l{\cel{3}instrumento, pleenda, en ensemblo. - DEFIS.}
\l{}
\l{partenero. La persono kun qua onu ludas kontre altra ludanti.}
\l{\cel{57}DEFR.}
}
